% Shared TikZ styles for the UML 2.5 use case diagrams of the report.
% Conventions followed (Ambler, The Elements of UML 2.0 Style; Microsoft UML use case
% guidelines): actors are roles drawn outside the system boundary, primary actors top-left,
% a specialised actor below its parent with the generalization arrow pointing to the parent
% (shared-target style: one hollow triangle under the parent, one trunk, one branch per
% child); use cases are horizontal ellipses named with a strong verb, stacked top to bottom
% in the order they happen; an included use case sits to the right of its base, an extending
% use case below its base; associations carry no arrowheads. All use cases are drawn alike;
% the implementation status is given in the report's status table, not in the diagram.
% Layout is placed by hand on a grid.
\usetikzlibrary{arrows.meta,shapes.geometric,positioning,calc,fit,backgrounds}

\definecolor{umlline}{HTML}{1F2937}

\tikzset{
  every node/.style={font=\sffamily\footnotesize},
  % a use case: an ellipse that hugs its (at most two-line) name
  uc/.style={ellipse, draw=umlline, line width=0.5pt, fill=white, align=center,
             inner xsep=0.3mm, inner ysep=0.2mm, minimum height=8mm, minimum width=26mm},
  assoc/.style={draw=umlline, line width=0.5pt},
  gen/.style={draw=umlline, line width=0.5pt, -{Triangle[open, length=2.6mm, width=2.6mm]}},
  dep/.style={draw=umlline, line width=0.5pt, dashed, -{Straight Barb[length=2mm, width=2.2mm]}},
  deplabel/.style={font=\sffamily\scriptsize, fill=white, inner sep=0.8pt},
  boundary/.style={draw=umlline, line width=0.7pt},
}

% \actor{node name}{position}{label}: a stick figure; the node is the figure's body.
\newcommand{\actor}[3]{%
  \begin{scope}[shift={#2}, line width=0.6pt, draw=umlline]
    \draw (0,0.34) circle (0.13);
    \draw (0,0.21) -- (0,-0.14);
    \draw (-0.23,0.08) -- (0.23,0.08);
    \draw (0,-0.14) -- (-0.18,-0.42);
    \draw (0,-0.14) -- (0.18,-0.42);
    \node[inner sep=0pt, minimum width=5mm, minimum height=10mm] (#1) at (0,0) {};
    \node[below=0.46cm of #1.center, align=center, font=\sffamily\footnotesize] (#1-label) {#3};
  \end{scope}%
}

% \specialize{child}{parent}{trunk x}: generalization from a child actor (placed to the
% right of the trunk) to its parent, drawn orthogonally: a horizontal branch from the child
% to the trunk, then up the trunk to the hollow triangle under the parent's name. Branches
% of several children share the trunk and the triangle (UML shared-target style).
\newcommand{\specialize}[3]{%
  \draw[gen] ($(#1.west)+(0.02,0)$) -- (#3,0 |- #1.west) -- ($(#2-label.south)+(0,-0.05)$);%
}
